\chapter{量子材料面试训练工作册}

\section{三级表达的使用方法}
每道重点题准备三个版本：
\begin{itemize}
\item \textbf{30 秒版：}定义或结论 $+$ 一条公式 $+$ 一句物理图像；
\item \textbf{2 分钟版：}补上推导骨架、适用条件、一个实验联系；
\item \textbf{黑板版：}从明确假设出发写关键中间步骤，并主动检查极限。
\end{itemize}
下面 30 道示范题给出三个版本的骨架。真正训练时请合上讲义，用自己的语言复述。

\newcommand{\trilevel}[4]{\subsection{#1}
\textbf{30 秒：}#2\par
\textbf{2 分钟：}#3\par
\textbf{黑板路线：}#4\par}

\trilevel{倒格子为什么不是纯数学技巧？}
{它是所有满足晶格周期相位条件的 Fourier 波矢集合，衍射动量转移和 Bloch 准动量都天然定义在倒空间。}
{由 $\vb a_i\cdot\vb b_j=2\pi\delta_{ij}$ 证明
$e^{i\vb G\cdot\vb R}=1$；晶面族对应法向倒格矢，第一布里渊区枚举不重复准动量。}
{构造 $\vb b_1=2\pi\vb a_2\times\vb a_3/\Omega_c$；证明相位条件；再画 Ewald/Bragg 几何。}

\trilevel{Bloch 定理的逻辑起点是什么？}
{周期势使 $[H,T_{\vb R}]=0$，平移本征值由群性质写成 $e^{i\vb k\cdot\vb R}$，因此 $\psi=e^{i\vb k\cdot\vb r}u_{\vb k}$。}
{说明离散平移、幺正本征值、$\vb k$ 模倒格矢等价，并指出准动量不等于机械动量。}
{从共同本征态出发逐步定义 $u_{\vb k}$，验证其周期性。}

\trilevel{区边界为什么开隙？}
{Bragg 条件使 $|\vb k\rangle$ 与 $|\vb k-\vb G\rangle$ 简并，周期势 $V_G$ 耦合二者，劈裂为 $2|V_G|$。}
{驻波在离子附近概率密度不同，所以势能不同；若对称性禁止矩阵元，交叉也可能保留。}
{写 $2\times2$ 简并微扰矩阵并对角化。}

\trilevel{低温声子热容为何是 $T^3$？}
{三维线性色散给 $g(\omega)\propto\omega^2$；Bose 能量积分给 $U\propto T^4$，所以 $C\propto T^3$。}
{说明 Debye 截止只固定总模数，低温上限可延至无穷；指数来自维数与色散。}
{由 $q$ 空间球计数推 $g(\omega)$，令 $x=\hbar\omega/\kb T$ 完成积分标度。}

\trilevel{电子热容为何线性？}
{Pauli 原理只允许费米面厚度 $\kb T$ 的电子激发，Sommerfeld 展开给 $C_V=\gamma T$。}
{参与粒子数 $\propto T$、每个激发能 $\propto T$，故 $\Delta U\propto T^2$；$\gamma$ 测态密度与质量重整化。}
{写 $U=\int Eg(E)f(E)\dd E$，给 Sommerfeld 二阶项和固定 $N$ 修正。}

\trilevel{Boltzmann 电导的核心是什么？}
{电场使平衡分布在费米面附近偏移，碰撞以 $\tau$ 恢复平衡；电流由 $v^2\tau(-\partial_Ef_0)$ 加权。}
{给出弛豫时间近似和张量电导，说明 Drude 是抛物线单带特例。}
{从稳态 Boltzmann 方程求 $\delta f$，代入 $j=-e\int v f$。}

\trilevel{费米面为什么比“所有占据态”重要？}
{低温热、磁与输运响应需要改变占据，深层态被 Pauli 阻塞，只有费米面附近可参与。}
{补充费米速度、曲率、拓扑和散射各向异性决定响应；费米海仍决定总密度等基态量。}
{用 $-\partial_Ef$ 在 $T\to0$ 变成 $\delta(E-E_F)$ 展示费米面投影。}

\trilevel{有效质量为什么可为负？}
{半经典加速度由能带曲率决定，$1/m^*=\hbar^{-2}\partial_k^2E$；带顶曲率负。}
{负质量电子等效为带正电、正质量空穴，适用于近带边波包。}
{联立 $\hbar\dot k=F$ 与 $v=\hbar^{-1}\partial_kE$。}

\trilevel{Mott 绝缘与能带绝缘有何本质差别？}
{能带绝缘来自满带和单粒子带隙；Mott 绝缘可在半填充带中由局域排斥阻止电荷移动。}
{用 Hubbard 的 $U/W$ 竞争说明电荷隙与低能自旋自由度，指出实际材料可能兼具晶场和关联。}
{写 Hubbard 模型，比较 $U=0$ 与 $U\gg t$ 的基态和激发。}

\trilevel{为什么 Hubbard 强耦合给反铁磁？}
{反平行自旋可通过双占据虚态获得二阶动能降低，平行自旋受 Pauli 阻止，故 $J=4t^2/U>0$。}
{强调虚过程不产生真实双占据，映射只在半填充、$U\gg t$ 的低能子空间成立。}
{对二聚体作二阶简并微扰，比较单态与三重态。}

\trilevel{交换相互作用为什么不是磁偶极作用？}
{它源于 Coulomb 相互作用结合费米波函数反对称性，能标通常远大于经典偶极作用。}
{区分直接交换、超交换、双交换和 RKKY；符号依轨道、占据与路径而定。}
{从两电子空间对称性和自旋对称性解释能量差，或从 Hubbard 推 $J$。}

\trilevel{Peierls 不稳定是什么？}
{一维金属的 $2k_F$ 晶格调制耦合两费米点并开隙，电子能降低可压过弹性能。}
{联系嵌套、Kohn 异常和电荷密度波，说明高维并非任意弱耦合都不稳定。}
{写二能级开隙矩阵，估算电子能与弹性能竞争。}

\trilevel{Berry 相位为何可观测？}
{Berry 连接虽依赖规范，闭合回路相位模 $2\pi$ 和 Berry 曲率是规范不变量，可影响干涉与半经典动力学。}
{联系绝热演化、反常速度、极化、量子振荡相位和拓扑数。}
{写 $\mathcal A=i\langle u|\nabla_k u\rangle$，演示规范变换及曲率不变。}

\trilevel{Chern 数为什么量子化？}
{它是整个二维布里渊区 Berry 曲率的 $2\pi$ 归一化通量，波函数规范拼接使其只能为整数。}
{整数在保持体能隙时不能连续改变，直接给满带 Hall 电导 $Ce^2/h$。}
{从二能带 $\hat d(\vb k)$ 映射到球面解释覆盖次数。}

\trilevel{拓扑相为何必须闭隙才能改变？}
{有隙时占据态投影可光滑连续变形，离散拓扑数不能变化；改变它必须出现奇点，通常是体隙关闭。}
{补充另一条路径是破坏保护对称性，使原分类不再适用。}
{用 Dirac 模型质量 $m$ 穿过零展示能隙 $2|m|$ 关闭和不变量改变。}

\trilevel{体边对应说了什么？}
{不同体拓扑不变量的界面不能处处保持同一有隙相，必须出现边界态或边界拓扑序。}
{用 SSH 端态、Chern 手性边缘和拓扑绝缘体 Dirac 表面态举例，并声明保护条件。}
{解质量变号 Dirac 方程的局域零模。}

\trilevel{Kramers 简并怎样保护表面态？}
{时间反演对半整数自旋满足 $\Theta^2=-1$，每个时间反演不变动量有正交同能伙伴。}
{非磁单粒子扰动不能把孤立 Kramers 对在该点混合开隙；磁性破坏时间反演后可开隙。}
{证明 $\langle\psi|\Theta\psi\rangle=0$，再讨论质量项的时间反演性质。}

\trilevel{Cooper 对为何能由任意弱吸引形成？}
{费米海阻塞了低能相空间，使费米面附近态密度有限，配对方程产生指数小但非零束缚能。}
{伙伴通常是时间反演态，BCS 中大量重叠对形成相干凝聚；“任意弱”限定于理想费米面和配对通道。}
{写 Cooper 方程并把求和变能量积分，展示对数发散。}

\trilevel{BCS 能隙的物理意义是什么？}
{它是 Bogoliubov 准粒子谱在费米面处的最低激发能标，反映电子空穴混合和配对序参量。}
{说明单粒子谱隙、破对能量与序参量在简单 $s$ 波模型中相关，但实验定义和多带体系需谨慎。}
{写 Nambu $2\times2$ 矩阵并求 $E_k=\sqrt{\xi_k^2+|\Delta|^2}$。}

\trilevel{Meissner 与零电阻有何区别？}
{零电阻是输运性质；Meissner 是平衡态排磁，说明超导不是普通完美导体。}
{用 London 穿透深度说明体内磁场衰减，完美导体则冻结初始磁通。}
{联立 London 与 Ampère 方程推 $\nabla^2B=B/\lambda_L^2$。}

\trilevel{非常规超导如何从实验识别？}
{关键是序参量对称性与节点，不是单纯 $T_c$ 高低；要结合低温幂律、相位敏感和谱学。}
{列比热、穿透深度、热导、NMR、Josephson 干涉及杂质效应，强调多证据交叉。}
{画 $s$ 波和 $d$ 波能隙随费米面角度变化，并推低能态密度。}

\trilevel{DFT 是多体理论还是单粒子理论？}
{基态密度泛函原则上包含相互作用，但 Kohn--Sham 把问题映射为辅助单粒子方程，近似集中在交换关联泛函。}
{解释自洽循环、能带解释边界、带隙低估和强关联失败。}
{写 Kohn--Sham 方程和密度更新，说明 $V_H,V_{xc}$。}

\trilevel{Wannier 函数为什么有用？}
{它把延展 Bloch 带变成局域轨道，构造紧束缚模型并高效插值能带、Berry 曲率和表面态。}
{说明规范选择、能带窗口和投影轨道会影响局域性，必须回验原始能带。}
{写 Bloch 到 Wannier 的 Fourier 变换。}

\trilevel{ARPES 实际测到什么？}
{近似测 $|M|^2A(\vb k,\omega)f(\omega)$，即占据单粒子谱函数乘矩阵元，而不是裸能带。}
{说明峰位置、宽度、强度的含义以及 $k_z$、表面敏感和分辨率限制。}
{由光电能量守恒和表面平行动量守恒写出 $E_B,k_\parallel$。}

\trilevel{STM 的 $\dd I/\dd V$ 为什么近似局域态密度？}
{弱隧穿且针尖态密度平坦时，电流是能窗内样品 LDOS 积分，微分选出 $E_F+eV$ 附近。}
{补充矩阵元、针尖结构、温度和锁相调制会卷积谱。}
{从 Fermi 黄金规则写隧穿电流的态密度卷积。}

\trilevel{量子振荡怎样测费米面？}
{Landau 能级周期穿过费米能，振荡频率由 Onsager 关系对应极值截面积。}
{温度衰减给回旋质量，场衰减给量子寿命，相位可含 Berry 信息但需处理 Zeeman 和三维修正。}
{写 $F=\hbar A_{\rm ext}/2\pi e$ 和 Lifshitz--Kosevich 因子。}

\trilevel{线性响应的统一图像是什么？}
{弱外场引起的响应是外场与因果响应函数的卷积，Kubo 公式把响应写成平衡对易关联。}
{静态极限给磁化率与压缩率，频率依赖给电导和介电函数，虚部表示耗散。}
{写 $\chi_{AB}(t)=i\theta(t)\langle[B(0),A(t)]\rangle/\hbar$。}

\trilevel{Landau 理论为何既强大又有限？}
{它只用序参量、对称性和解析展开就给相结构与平均场指数，但忽略强临界涨落和微观机制。}
{说明梯度项、关联长度、上临界维数以及一级相变可能来源。}
{最小化 $a(T-T_c)m^2+bm^4-hm$，求 $m,\chi$。}

\trilevel{怎样区分经典相变和量子相变？}
{经典相变由温度驱动、热涨落主导；量子相变发生在零温，由非热参数调控、量子涨落主导。}
{有限温度量子临界扇区由 $\hbar/\kb T$ 截断虚时间，出现动态指数 $z$。}
{画参数--温度相图，标出有序相、量子临界点和交叉区。}

\trilevel{给一张陌生数据图时怎样开始？}
{先读坐标、单位、控制变量和误差；再识别极限与特征尺度；最后列最简解释和至少一个替代解释。}
{检查原始量还是处理量、是否归一化、是否跨样品比较，再把特征对应到模型可观测量。}
{现场重画示意图，标出斜率、截距、峰位、宽度、积分权重和不确定性。}

\section{六十道核心训练题}
\begin{multicols}{2}
\begin{enumerate}
\item Wigner--Seitz 原胞为何保留点群对称？
\item 晶面指数与倒格矢如何对应？
\item Ewald 球怎样表示弹性衍射？
\item diamond 消光规律怎样得到？
\item 声学和规则是什么？
\item LO--TO 劈裂来自哪里？
\item Einstein 模型为何低温失败？
\item 声子准动量何时守恒？
\item Normal 与 Umklapp 过程区别？
\item Debye 温度的物理含义？
\item Sommerfeld 模型比 Drude 多了什么？
\item Pauli 顺磁为何弱温度依赖？
\item Landau 抗磁为何是量子效应？
\item Van Hove 奇点如何产生？
\item 为什么满带不导电？
\item 波包群速度怎样从能带得到？
\item 空穴的能量和动量如何定义？
\item 直接/间接带隙怎样用光学区分？
\item PN 结耗尽层怎样形成？
\item Wiedemann--Franz 何时失效？
\item Hall 符号为何不能总判载流子类型？
\item 磁阻可能来自哪些机制？
\item RKKY 相互作用为何振荡？
\item 铁磁与反铁磁自旋波有何差别？
\item Mermin--Wagner 定理的直觉是什么？
\item 顺磁、铁磁和自旋玻璃如何区分？
\item 密度波的序参量为什么可为复数？
\item Kohn 异常是什么？
\item Fermi 液体准粒子是什么？
\item 自能实部和虚部分别改变什么？
\item 屏蔽为何依赖维数和态密度？
\item 等离激元为何是集体模？
\item 反常 Hall 与 Berry 曲率如何联系？
\item 半经典方程中反常速度是什么？
\item Zak 相位与极化如何联系？
\item 对称性指标能解决什么问题？
\item 反演本征值怎样帮助判断拓扑？
\item Weyl 点与 Dirac 点有何区别？
\item Fermi arc 是什么？
\item 拓扑保护不等于绝对不散射，为什么？
\item Andreev 反射是什么？
\item Josephson 两条关系是什么？
\item I 型与 II 型超导区别？
\item 相干长度与穿透深度有什么意义？
\item BCS--BEC crossover 的直觉是什么？
\item 多带超导会带来什么现象？
\item 为什么零偏压峰不是 Majorana 的充分证据？
\item Kohn--Sham 能隙为何常被低估？
\item DFT+$U$ 的 $U$ 是什么？
\item 自洽计算怎样判断真正收敛？
\item XRD 峰宽能反映什么？
\item 中子为何适合测磁激发？
\item ARPES 如何区分表面态与体态？
\item STS 相干峰说明什么？
\item 量子振荡为何要求高迁移率？
\item 比热跳变如何估算超导体积分数？
\item Curie 尾如何识别磁杂质？
\item 光学谱权重和规则是什么？
\item 泵浦--探测的时间尺度如何解释？
\item 科研报告中怎样诚实表达不确定性？
\end{enumerate}
\end{multicols}

\section{三十道追问与科研衔接题}
\begin{enumerate}
\item 你的研究对象为何可以用单粒子能带描述？何时会失效？
\item 你采用的边界条件会不会改变结论？
\item 若去掉自旋轨道耦合，结果会怎样？
\item 若破坏时间反演或反演，对简并有什么影响？
\item 你模型中的能标大小如何从实验估计？
\item 你的近似小参数是什么？
\item 如何检查数值结果不是有限尺寸伪影？
\item 网格、截断和展宽参数怎样做收敛测试？
\item 为什么选择这个基底或规范？
\item 结果是否满足粒子数、谱权重或和规则？
\item 如何用量纲和极限检查代码？
\item 若实验与 DFT 不符，可能有哪些原因？
\item 如何区分关联重整化与普通带宽误差？
\item 如何区分体态、表面态与杂质态？
\item 如何判断观察到的是相变还是 crossover？
\item 序参量是什么，破坏了什么对称性？
\item 有哪些竞争序，怎样实验区分？
\item 为什么一张拟合曲线不能证明机制？
\item 参数之间是否高度相关，误差如何估计？
\item 你的实验/计算直接测量或输出的量是什么？
\item 从原始数据到结论经过了哪些处理？
\item 背景扣除是否会制造峰或能隙？
\item 替代模型能否同样解释数据？
\item 下一项最有判别力的实验是什么？
\item 若样品存在无序，结论如何变化？
\item 温度、磁场、压力、掺杂中哪个是最自然的调控参量？
\item 如何把你的课题放到量子材料的大图景中？
\item 本科项目中你真正独立完成了哪一部分？
\item 如果结果与预期相反，你如何排查？
\item 进入课题组后前三个月你准备补什么能力？
\end{enumerate}

\section{不会时的安全回答}
\begin{interviewbox}
推荐句式：“我目前不记得完整系数，但我可以先从对称性和量纲判断其形式。若采用……近似，起点应是……；接下来需要计算……才能确定。这个近似在……时会失效。”

避免三种做法：把猜测说成事实；用更多术语掩盖逻辑空缺；在面试官纠正后仍坚持原表述。承认边界、继续推理，往往比背错结论更能展示研究潜力。
\end{interviewbox}

